\subsection{Values and Expressions}
In ordinary mathematics, we usually write operators between their operands:
$2+3$, $6/2$, or $x^2$.  This is called \term{infix notation}.  Racket
instead uses a uniform notation for arithmetic and function application.

\subsubsection{Beginning Student source files}
In DrRacket, select the \term{Beginning Student} language.  When writing a
Racket source file, the corresponding language declaration is written at the
beginning of the file:
\begin{minted}{text}
#lang htdp/bsl
\end{minted}
The declaration chooses the language; it is not an expression that evaluates
to a number or Boolean value.  The expressions and definitions come after it:
\begin{minted}{text}
#lang htdp/bsl

(define x 7)
(define y 11)
(+ x y)       ; evaluates to 18
\end{minted}
The line beginning with a semicolon is a comment.  In the examples below, the
language declaration is omitted when only an individual expression is shown.

\begin{definition}[Prefix notation]
A Racket expression is fully parenthesized, and its operator is written first:
\begin{minted}{text}
(+ 2 3)
(/ 6 2)
(sqr x)
(expt x 2)
\end{minted}
The expression after the operator is its argument list.  The parentheses make
the grouping explicit, so we do not need a separate precedence convention.
The operator-first idea is also called \term{Polish notation}.  For example,
classical Polish notation writes $2+3$ as \texttt{+ 2 3}, without parentheses.
Racket uses a parenthesized form of this idea, \texttt{(+ 2 3)}, so the
parentheses explicitly show where the expression begins and ends.  Thus it is
helpful to think of Racket's prefix notation as parenthesized Polish notation,
while using \term{prefix notation} as the standard CS135 term.

For example, the mathematical expressions
\[
  8-2+3 \qquad\text{and}\qquad 8-(2\mathbin{\cdot}3)
\]
become, respectively,
\begin{minted}{text}
(+ (- 8 2) 3)
(- 8 (* 2 3))
\end{minted}
\end{definition}

\subsubsection{Evaluating an expression}
Evaluation is a process of substitution.  Each step replaces a subexpression
by its value, and the result is still a valid (but simpler) Racket program.  A
substitution step chooses the leftmost, innermost subexpression with no
parentheses inside it.  We write the yields symbol $\Rightarrow$ between steps.

\begin{example}
\begin{align*}
&(* (/ 6 2) (+ 1 2)) \\
&\Rightarrow (* 3 (+ 1 2)) \\
&\Rightarrow (* 3 3) \\
&\Rightarrow 9.
\end{align*}
\end{example}

\subsubsection{Basic arithmetic operators}
The arithmetic operators $+$, $-$, $*$, and $/$ accept two or more values in
a single step (except that $-$ also supports one argument for negation).
\begin{align*}
  &( +\ 11\ 12\ 13) \Rightarrow 36, \\
  &( /\ 720\ 5\ 4\ 3\ 2\ 1) \Rightarrow 6, \\
  &(-\ 144) \Rightarrow -144, \qquad (-\ {-21}) \Rightarrow 21.
\end{align*}
In actual Racket syntax, spaces separate the arguments:
\begin{minted}{text}
(+ 11 12 13)
(/ 720 5 4 3 2 1)
(- 144)
(- -21)
\end{minted}
Thus, for example, $(+\ 72)$ is an error: addition requires at least two
arguments.

\subsection{Exact numbers}
The first number systems in CS135 are exact.  Exactness means that the value
is represented without rounding; for example, $3/2$ is not silently replaced
by a decimal approximation.

\begin{definition}[Natural numbers]
The natural numbers are defined recursively by
\[
  0\in\N, \qquad n\in\N\Longrightarrow n+1\in\N.
\]
CS135 writes the type as \texttt{Nat} in Racket discussions.  Starting at
zero is especially convenient for computer science.
\end{definition}

\begin{definition}[Integers]
The integers are
\[
  \Z=\{\ldots,-3,-2,-1,0,1,2,3,\ldots\}.
\]
CS135 writes \texttt{Int} for this type, which extends \texttt{Nat} by adding
negative numbers.
\end{definition}

\begin{definition}[Rational numbers]
The rational numbers are numbers of the form $p/q$, where $p,q\in\Z$ and
$q\ne0$.  CS135 writes \texttt{Rat}.  Racket represents rational numbers
exactly: \texttt{(/ 3 2)} evaluates to the exact number $3/2$.  Here
\texttt{(/ 3 2)} is an expression performing division, while \texttt{3/2}
is a literal way to write the resulting number.
\end{definition}

In Racket, exact \texttt{Nat}, \texttt{Int}, and \texttt{Rat} values can be
arbitrarily large or small, limited only by available memory.  ``Arbitrarily''
does not mean infinitely large: every physical computer has finite memory.

\subsubsection{Integer division and type relationships}
The operators \texttt{quotient} and \texttt{remainder} perform integer
division without producing a rational number:
\begin{minted}{text}
(quotient 17 5)  ; 3
(remainder 17 5) ; 2
\end{minted}
They consume two \texttt{Int}s and produce an \texttt{Int}; applying them to
rational values is an error.

The types form a subtype chain:
\[
  \texttt{Nat}\subseteq\texttt{Int}\subseteq\texttt{Rat}.
\]
Thus every natural number is an integer, and every integer is rational.  If an
operator works for a type, it automatically works for its subtypes; an
operator may still have additional restrictions, as \texttt{remainder} does.
The subtype idea will recur throughout computer science.

\subsubsection{Other numerical operators}
The operators \texttt{max} and \texttt{min} select the largest and smallest
argument:
\begin{align*}
  &\texttt{(max 4 8 -9 3)} \Rightarrow 8, &
  &\texttt{(min 4 8 -9 3)} \Rightarrow -9.
\end{align*}
Unlike $+$, $*$, and $/$, each also accepts a single argument:
\[
  \texttt{(max 9)}\Rightarrow9, \qquad
  \texttt{(min 9)}\Rightarrow9.
\]
Exponentiation is written with \texttt{sqr} and \texttt{expt}:
\[
  \texttt{(sqr 12)}\Rightarrow 12^2=144, \qquad
  \texttt{(expt 3 4)}\Rightarrow 3^4=81.
\]

\subsubsection{Real numbers}
The rationals are enough for addition, subtraction, multiplication, division,
and integer powers.  But $\sqrt{2}=2^{1/2}$ is irrational, so we need the real
numbers $\R$.  An irrational real number cannot be represented exactly in
finite memory.  We return to this distinction between exact and inexact
numbers in Lecture 02.

\subsection{Constant definitions}
A mathematical variable may change as we work, but a Racket name introduced
with \texttt{define} is a \term{constant}: it is defined once and does not
change.
\begin{minted}{text}
(define x 7)
(define y 11)
\end{minted}
The first definition binds the value $7$ to the name \texttt{x}; the second
binds $11$ to \texttt{y}.  We can then write, for example,
\begin{minted}{text}
(+ (sqr x) (* 6 x y) (sqr y) (* 9 x) (- (* 3 y)) -100)
\end{minted}

A definition is evaluated by substitution.  For instance,
\begin{minted}{text}
(define y (+ x 4))
\end{minted}
first substitutes the previously defined value of \texttt{x}, giving
\texttt{(+ 7 4)}, and then evaluates that expression to $11$ before binding
$11$ to \texttt{y}.  A name must be defined before it is used, and a name
cannot be redefined once it has been defined.

Meaningful constants also prevent \term{magic numbers}, whose purpose is
otherwise hidden in a program:
\begin{minted}{text}
(define lightspeed 299792458)
(define hours-in-day 24)
(define feet-in-mile 5280)
(define life-the-universe-and-everything 42)
\end{minted}

\subsection{Boolean expressions}
In mathematics, the statement $x<5$ describes a condition on a possibly
unknown $x$.  In Racket, a constant such as \texttt{x} already has a value, so
\texttt{(< x 5)} is an expression asking whether that value is less than $5$.
For example:
\begin{minted}{text}
(define x 4)
(< x 6)  ; true
(> x 6)  ; false
(= x 7)  ; false
(>= 5 x) ; true
(<= 5 x) ; false
\end{minted}
The comparison operators \texttt{<}, \texttt{>}, \texttt{<=}, \texttt{>=}, and
\texttt{=} consume numbers and produce a Boolean value.  The type is
\texttt{Bool}, whose only values are \texttt{true} and \texttt{false}.

\subsubsection{Combining Boolean values}
The operators \texttt{and}, \texttt{or}, and \texttt{not} combine Boolean
expressions:
\begin{itemize}
  \item \texttt{and} is true exactly when all its arguments are true;
  \item \texttt{or} is true when at least one argument is true;
  \item \texttt{not} reverses a Boolean value.
\end{itemize}
Both \texttt{and} and \texttt{or} require at least two arguments, whereas
\texttt{not} takes one.  For example,
\begin{align*}
&\texttt{(and (or (< 2 3) (> 1 5)) (not (< 10 15)) (<= 5 5))} \\
&\Rightarrow \texttt{(and (or true (> 1 5)) (not (< 10 15)) (<= 5 5))} \\
&\Rightarrow \texttt{(and true (not true) (<= 5 5))} \\
&\Rightarrow \texttt{(and true false (<= 5 5))} \\
&\Rightarrow \texttt{false}.
\end{align*}

\subsubsection{Short-circuit evaluation}
The definitions of \texttt{and} and \texttt{or} let Racket stop early.  If
one argument of \texttt{and} is false, the whole expression is false; if one
argument of \texttt{or} is true, the whole expression is true.  This is called
\term{short-circuit evaluation}:
\begin{align*}
  &\texttt{(or (< 2 3) (> 1 5))}
     \Rightarrow \texttt{(or true (> 1 5))}\Rightarrow\texttt{true},\\
  &\texttt{(and true (not true) (<= 5 5))}
     \Rightarrow \texttt{(and true false (<= 5 5))}\Rightarrow\texttt{false}.
\end{align*}

\subsubsection{De Morgan's laws}
For Boolean expressions $A$ and $B$, the familiar set-theoretic identities
are also valid:
\begin{align*}
\texttt{(not (and A B))}
  &\equiv \texttt{(or (not A) (not B))},\\
\texttt{(not (or A B))}
  &\equiv \texttt{(and (not A) (not B))}.
\end{align*}
For example,
\begin{align*}
&\texttt{(not (or (< 2 3) (> 1 5)))} \\
&\equiv \texttt{(and (not (< 2 3)) (not (> 1 5)))} \\
&\equiv \texttt{(and (>= 2 3) (<= 1 5))}.
\end{align*}

\subsection{Lecture 01 summary}
Lecture 01 builds the first layer of a mathematical model of computation:
values and expressions, then constants and Boolean expressions.  The next
lectures add functions, conditional expressions, and recursion; together these
will form a complete small model of a computer.  Lists will later make it
possible to organize data and study algorithms and data structures.

The main ideas to remember are:
\begin{itemize}
  \item write mathematical operations in fully parenthesized prefix notation;
  \item evaluate by valid substitution steps until a value is reached;
  \item distinguish exact \texttt{Nat}, \texttt{Int}, and \texttt{Rat} values;
  \item treat \texttt{define}d names as constants; and
  \item use Boolean expressions to represent true-or-false questions.
\end{itemize}

\subsubsection{Allowed constructs}
The Lecture 01 allowed-constructs list is:
\begin{minted}{text}
( ) + - * / = < > <= >= and define expt false max min not or
quotient remainder sqr true Bool Int Nat Rat
\end{minted}
Assignments may permit the constructs from a particular lecture, together
with any restrictions stated on the assignment's first page.
